% =====================================================================
%  29.5  The method boundary and a dominator preview.
%  Exports: sec:c31-boundary  def:c31-signature  def:c31-pick
%  fig:c31-taxonomy2
%  WORKED section only.  Decides WHEN a graph has a useful tree inside it (and
%  which of the four decompositions) and when it does not; hands off to Ch.31
%  (MST), Ch.30 (counting labelled trees), Ch.26 (virtual trees) and graph-land
%  (flows, dominators).  A one-paragraph dominator preview closes the chapter
%  with the honest caveat that flow-graph reachability is not a tree-of-
%  components story.
% =====================================================================
\section{The Method Boundary: Which Decomposition}\label{sec:c31-boundary}

\begin{previouslybox}[Four decompositions, one question each]
\(\secref{c31-bridges}\) kills unavoidable \emph{edges} (bridge tree);
\(\secref{c31-block}\) kills hinge \emph{vertices} (block-cut tree);
\(\secref{c31-functional}\) reads an out-degree-one map (rho pieces);
\(\secref{c31-cactus}\) exploits the case where every block is a lone cycle.  The
boundary job is a \emph{decision}: given a fresh statement, name the redundancy to
remove -- or admit the graph has no tree worth extracting.
\end{previouslybox}

\begin{topicintro}
\textbf{Prerequisites.} All of 29.1--29.4.  This section adds no algorithm; it adds a
\emph{signature} you run a statement through and a list of the neighbours that own
what this chapter rejects.

\textbf{After this section you can}
\begin{itemize}[nosep]
  \item state the four-question signature that identifies a graph-to-tree job; \hfill\lv{4}
  \item pick bridge-tree / block-cut / functional / cactus from a statement; \hfill\lv{5}
  \item hand a mis-fit problem to the right chapter (MST, counting, virtual trees, flow); \hfill\lv{5}
  \item explain, honestly, what a dominator tree is and why it is mostly \emph{not} here. \hfill\lv{5}
\end{itemize}
\end{topicintro}

% ---------------------------------------------------------------
\subsection{The recognition signature}
% ---------------------------------------------------------------

\begin{definition}{The graph-to-tree signature (four tests)}{c31-signature}
Run a statement through these; you are in this chapter when all four point the same way.
\emph{(T1) Does the answer depend on connectivity that survives deleting a fixed edge or
vertex?} If yes, an edge/vertex decomposition (bridge tree or block-cut) is on the table.
\emph{(T2) Is there exactly one outgoing choice per object} (a map, a vote, a single
pointer)? If yes, it is a functional graph \(\to\) rho.
\emph{(T3) Are cycles guaranteed to be edge-disjoint} (``each edge in at most one cycle'',
\(m\approx n\))? If yes, cactus.
\emph{(T4) Once the redundant part is contracted, is the remaining object a tree you can
run Ch.~12/13 on}? If yes, the extraction pays; if no (the residual still has interacting
cycles), you are leaving the chapter.
\end{definition}

\begin{recognitionbox}[Which of the four am I?]
\reckey{``removing one edge/vertex disconnects'', ``on exactly one simple cycle'',
``each \(x\) maps to one \(f(x)\)'', ``eventually cycles'', ``each edge in at most one
cycle'', ``is it a cactus / tree'', ``orient so you can always return''}
\recuse{
  \item \textbf{Edge-avoidance / orientation} questions \(\to\) bridge tree
        (\(\secref{c31-bridges}\)); ``is it bridgeless'' even decides strong orientation
        (Robbins, \(\thmref{c31-robbins}\)).
  \item \textbf{Vertex-avoidance / shared-hinge} questions \(\to\) block-cut tree
        (\(\secref{c31-block}\)); lone-cycle edge queries live here (CF 962F).
  \item \textbf{A single arrow per vertex} \(\to\) functional rho (\(\secref{c31-functional}\)).
  \item \textbf{Edge-disjoint cycles} \(\to\) cactus contraction + ring transfer
        (\(\secref{c31-cactus}\)).
}
\recnotuse{
  \item \textbf{Cheapest total connection / add edges to connect} \(\to\) minimum
        spanning tree: Ch.~31 (\chref{c17}).
  \item \textbf{How many labelled trees satisfy \(\ldots\)} \(\to\) counting:
        Ch.~30 (\chref{c32}); the graph is not given, the family of trees is.
  \item \textbf{Steiner tree over a few marked vertices of a weighted graph} \(\to\)
        virtual trees: Ch.~26 (\chref{c30}).
}
One-line test: \emph{can I name a redundancy (a bridge, a hinge, a self-loop of a map,
a cycle block) whose removal leaves a tree, and does the answer only care about the
tree?} If not, this chapter is the wrong tool.
\end{recognitionbox}

\subsection{The pick table}

\begin{definition}{The decomposition pick}{c31-pick}
Given the signature verdict (\(\defref{c31-signature}\)), the table below fixes the
object to build and the black box to run on it: robustness-to-edge-deletion uses the
bridge tree, robustness-to-vertex-deletion the block-cut tree, a single-arrow map the
functional rho, and edge-disjoint cycles the cactus contraction.
\end{definition}

\begin{center}
\renewcommand{\arraystretch}{1.25}
\setlength{\tabcolsep}{3pt}
\begin{tabular}{@{}p{0.30\linewidth} p{0.24\linewidth} p{0.20\linewidth} >{\raggedright\arraybackslash}p{0.16\linewidth}@{}}
\toprule
\textbf{If the query is\dots} & \textbf{Build} & \textbf{Then run} & \textbf{Section}\\
\midrule
``must every path use this edge?'' & bridge tree & Ch.13 DP on it & \(\secref{c31-bridges}\)\\
``which edges sit on one cycle?'' & block-cut tree & \(|E|=|V|\) test & \(\secref{c31-block}\)\\
``after how many steps in a cycle?'' & functional rho & \(\rho,\mu\) DP & \(\secref{c31-functional}\)\\
``count X over the whole graph'' & cactus block tree & ring transfer + fold & \(\secref{c31-cactus}\)\\
\bottomrule
\end{tabular}
\end{center}

\subsection{A dominator preview (and where the full story lives)}

\begin{notebox}[One paragraph on dominators -- then we bow out]
Fix a start vertex \(s\) in a \emph{directed} graph.  Vertex \(d\) \textbf{dominates}
\(v\) if every path \(s\leadsto v\) passes through \(d\); the relation is a tree (the
\term{dominator tree}) rooted at \(s\).  It is the directed, vertex-level cousin of
this chapter's edge/vertex decompositions, and there is a genuinely pretty bridge of
ideas: in an undirected graph, \(u\!-\!v\) is a bridge \emph{iff} \(v\) dominates
\(u\) in the DFS from either endpoint \emph{and} \(u\) dominates \(v\) from the other
-- bridges are two-directional domination.  But the general (directed) dominator tree
needs the Lengauer--Tarjan algorithm, and dominators answer \emph{flow reachability},
not ``contract the redundancy and DP''; the residual object is not the tree-of-
components this chapter builds.  That is graph-algorithm territory, deliberately out
of scope here: we preview the analogy so you recognise a dominator problem and go
read about it properly rather than force a bridge tree onto it.
\end{notebox}

\begin{comparisonbox}[This chapter vs.\ its neighbours]
\emph{Bridge/block-cut tree (here).} A tree extracted from \emph{one given graph} to
answer connectivity-robustness questions about that graph.
\emph{MST (Ch.~31, \chref{c17}).} A tree chosen from weighted edges to \emph{optimise
total cost}; no decomposition, a greedy exchange argument.
\emph{Counting (Ch.~30, \chref{c32}).} No graph given at all -- count how many trees
meet a specification (Prüfer, Matrix-Tree).
\emph{Virtual tree (Ch.~26, \chref{c30}).} A tree spanning marked vertices \emph{of an
existing tree} via LCA -- compression, not decomposition.  Four different ``trees
inside graphs''; only the first two are decomposition, and only the first is this
chapter's.
\end{comparisonbox}

\begin{stuckbox}[If no decomposition seems to fit]
\begin{itemize}[nosep,leftmargin=1.4em]
  \item Are you sure the object is a single given graph? If it is ``how many graphs/trees
        satisfy \(\ldots\)'', you want Ch.~30 (\chref{c32}).
  \item Is there a weight you must minimise? Then it is likely an MST problem
        (Ch.~31, \chref{c17}), not a robustness problem.
  \item Is the graph \emph{directed} and the question about reaching / cutting all
        paths? Strong-connectivity condensation is the directed analogue of the bridge
        tree; dominators (\(\secref{c31-boundary}\)'s preview) are the harder cousin.
  \item Do cycles share edges all over the place (the graph is ``tree-like'' but not a
        cactus)? Then the block-cut tree exists but each block is dense -- you are at
        the boundary; a bounded cyclomatic number is the only rescue
        (\(\secref{c31-cactus}\)'s hardening note).
\end{itemize}
\end{stuckbox}

\begin{warningbox}[Three ways to leave the chapter wrongly]
\begin{enumerate}[nosep,leftmargin=1.6em]
  \item \textbf{Confusing the bridge tree and the block-cut tree.} Edge-robustness uses
        2-edge-components (\(\secref{c31-bridges}\)); vertex-robustness uses blocks
        (\(\secref{c31-block}\)).  They differ on exactly the ``cycles sharing a vertex''
        case -- pick by \emph{what you delete}, an edge or a vertex.
  \item \textbf{Forcing a cactus DP on a non-cactus.} If a block has \(|E|>|V|\) the ring
        transfer is invalid; detect it in the block pop and stop (\(\secref{c31-cactus}\)).
  \item \textbf{Solving a counting problem with a decomposition.} ``Number of spanning
        trees of this graph'' is the Matrix-Tree theorem (Ch.~30, \chref{c32}), not a
        bridge/block question -- no edge is special there.
\end{enumerate}
\end{warningbox}

\begin{recallbox}[29.5 in one breath]
\emph{The signature is four tests -- edge-robustness, vertex-robustness, single-arrow,
edge-disjoint cycles -- and it hands off everything else: weights to MST (\chref{c17}),
enumeration to counting (\chref{c32}), marked-vertex compression to virtual trees
(\chref{c30}), directed flow cuts to dominators (preview only, out of scope).  Choose the
redundancy you can delete, contract it, and reuse Ch.~12/13 on the tree that remains
(\(\defref{c31-signature}\), \(\defref{c31-pick}\)).}
\selftest{
  \item A query asks ``does deleting vertex \(c\) disconnect the graph?'' -- which decomposition? \emph{Block-cut tree: cut vertices are the hinges (\(\secref{c31-block}\)); a bridge tree answers edge deletions, not vertex. \(\defref{c31-signature}\).)}
  \item The graph has weights and you must connect all vertices cheapest -- still this chapter? \emph{No -- that is a minimum spanning tree, Ch.~31 (\chref{c17}); T4 fails because there is no redundancy to contract, you choose edges by cost. \(\defref{c31-pick}\).)}
  \item In one sentence, why is the dominator tree mostly out of scope? \emph{It answers directed flow-reachability (every \(s\)-path through \(d\)); the residual object is not a tree-of-components and needs Lengauer--Tarjan, not a bridge/block decomposition.}
}
\end{recallbox}
